%% ma: L12-B08-0091
%% lop: 12
%% chuong: 2
%% bai: 8
%% dang: 12.8.12
%% loai: TN4
%% muc: TH
%% dap_an: "D"
%% nguon: "CĐ2.2.pdf giáo viên gửi 27/10/2026 (M07, Lê Minh Khôi – Nguyễn Anh Khôi, Chương II Chủ đề 2 Tọa độ của vectơ trong không gian), Câu 25"
%% kien_thuc: "Bài 8 (tọa độ vectơ, hai vectơ bằng nhau)"
%% trang_thai: chinh-thuc
\begin{ex}
Trong không gian $Oxyz$, cho ba điểm $M(1;-1;1)$, $E(2;0;-1)$ và $X(-1;2;1)$. Xét điểm $K(x_K;y_K;z_K)$ sao cho tứ giác $XKEM$ là hình bình hành. Tính giá trị của biểu thức $T=26y_K+27x_K-z_K$.
\choice{$T=80$}{$T=82$}{$T=77$}{\True $T=79$}
\loigiai{$XKEM$ là hình bình hành $\Leftrightarrow\overrightarrow{XK}=\overrightarrow{ME}=(1;1;-2)\Rightarrow K(0;3;-1)$.

$T=26\cdot3+0+1=79$.}
\end{ex}
